\section{Network Model and Assumptions}
\label{sec:model}

We model a quantum network as an undirected graph $G = (V, E)$ where nodes are repeater stations and edges are fiber links with known lengths. We consider three topologies: NSFNET (14 nodes, 21 edges, US backbone), a pruned 17-node SURFnet (Dutch research network, following Vardoyan et al.~\cite{vardoyan2024bipartite} Fig.~13 with real fiber distances from~\cite{pouryousef2024resource}), and Abilene/Internet2 (11 nodes, 14 edges, US backbone). Since NSFNET and Abilene have continental-scale distances (250--2300~km), we scale them down by $100\times$ to bring links into the quantum-viable regime (2.5--23~km). SURFnet's distances (15--78~km) are already short enough and used unscaled. This scaling approach is standard in the literature~\cite{vardoyan2024bipartite, pouryousef2024resource}.

For each topology we pick three source-destination pairs spanning opposite ends of the network: (WA, DC), (CA1, NY), (TX, MI) for NSFNET; (AMS, NIJ), (DEL, ENS), (RTM, ZWO) for SURFnet; and (SEA, NYC), (LAX, CHI), (HOU, WAS) for Abilene. All three pairs compete for shared resources simultaneously at every time step.

The simulation runs in discrete time steps, each consisting of four phases: link generation, decoherence, routing, and swapping/delivery. Entanglement generation follows a Bernoulli process: every edge independently attempts to generate one Bell pair per time step, succeeding with probability $p_\text{gen} = p_0 \cdot \exp(-L_e / L_\text{att})$, where $p_0 = 1.0$ and $L_\text{att} = 22$~km. We assume one communication qubit per link interface, so all edges attempt generation in parallel. Successful pairs are stored with initial fidelity $F_\text{init} = 0.95$ in quantum memories at both endpoints. Memories use a deterministic cutoff: links older than $T_\text{cut} = 10$ time steps are discarded. We do not model continuous fidelity decay during storage.

Entanglement swapping at intermediate repeaters uses the depolarizing noise model: $F_\text{out} = F_1 F_2 + (1 - F_1)(1 - F_2)/3$, applied sequentially along the path. Swaps are deterministic ($q_\text{swap} = 1.0$) and gates are assumed perfect (no noise beyond the depolarizing model). A delivered pair counts as successful only if its end-to-end fidelity meets a threshold $F_\text{min} = 0.80$. With $F_\text{init} = 0.95$, this means paths longer than 4 hops cannot produce usable entanglement.

We assume free classical communication: all nodes have instantaneous global knowledge of generation outcomes at every time step. This lets routing algorithms observe the full entanglement state before making decisions. In a real network, classical signaling delays would force decisions based on partial or stale information, hurting MCF (which relies on global state) more than greedy routing. Our results therefore represent an upper bound on both algorithms' performance, with the MCF advantage likely narrowing under realistic delays.

For reference, a summary of all assumptions: (1) Bernoulli generation with $p_\text{gen} = p_0 \cdot e^{-L/L_\text{att}}$; (2) one qubit per link interface (parallel generation); (3) deterministic memory cutoff, no fidelity decay in storage; (4) depolarizing swap model with deterministic success; (5) perfect local gates; (6) free classical communication; (7) fidelity threshold $F_\text{min}$ for delivery; (8) NSFNET and Abilene scaled $\div 100$, SURFnet unscaled; (9) three concurrent user pairs per topology chosen across opposite ends of the network.
